\scanpage{27}
\mspage{1}
$X$ surface \deleted{conforme}\inserted{hol \unsure{0}-connexe} (\unsure{pas néc cpte})

$\widetilde X$ rev. universel

\margin{\uline{NB} Ici ``surface'' sous-entend ``séparée''}

\uline{Th} Conditions équivalentes

(i) $X$ n'est pas isom. à une des trois surfaces
\[
\mathbb P^1_{\mathbb C},\ \mathbb E^1_{\mathbb C}=\mathbb P^1_{\mathbb C}\backslash\{\infty\},\quad G_{m\,\mathbb C}=\mathbb P^1_{\mathbb C}\backslash\{\infty,0\},
\]
ni à une courbe elliptique complexe

(ii) $\widetilde X$ est isomorphe à $H$ (demi-plan de Poincaré)

(iii) Toute application holomorphe $\mathbb C\longrightarrow X$ est constante

(iv) [Si $X$ est un ouvert de une droite projective complexe\deleted{]} (resp. d'une droite affine complexe)]
\[
\operatorname{card}(\deleted{\unsure{\mathbb P^1}}X'\backslash X)\ge3
\quad(\text{resp. }\operatorname{card}(\deleted{\unsure{\mathbb A^1}}X'\backslash X)\ge2).
\]

\deleted{(iv') [Si $X$ est un ouvert de \unsure{C. unif.}}

(v) [Si $X=\widehat X\backslash S$, $\widehat X$ compacte\inserted{genre $g$}, $S$ finie\inserted{card $\nu$}]
\deleted{holomorphe connexe $X'$}
\[
\unsure{2|3}g+\nu\ge3\qquad\text{i.e. }(g,\nu)\unsure{\notin|\ne}\{(0,0)(0,1)(0,2)(1,0)\}
\]
\[
\text{i.e. }\chi_c(X)<0
\]

\uline{Déf} On dit alors que $X$ est hyperbolique

\uline{Corollaire 1} Un ouvert \unsure{0}-connexe d'une surface hol. hyperbolique est hyperbolique

\uline{Corollaire 2} Soit $X$ surface hol. \uline{à bord}, \deleted{\unsure{0-co}} connexe, avec $\partial X\ne\varnothing$. Alors $\operatorname{Int}X$ est hyperbolique
% archive foliation: 29
